\chapter{Conclusioni e Ricerca
Futura}\label{conclusioni-e-ricerca-futura}

\begin{quote}
\textbf{Argumento central:} la DEQ$^2$ è una teoria \emph{aperta} ---
esposta alla falsificação (Cap.~11), strutturata per estensione (paper
EN, articolo PT, dataset multi-paese), inserita esplicitamente in una
tradizione (Arrighi, Turchin, Wallerstein, Cervo, Marcuse-Adorno, Taleb)
di cui rivendica l'eredità senza pretendere di esaurirla. Questo
capítulo conclude il libro come \emph{atto epistemico e politico}:
chiude la trattazione, dichiara cosa resta da fare, invita
esplicitamente al dissenso informato.
\end{quote}



\section{Riepilogo della Tesi
Centrale}\label{riepilogo-della-tesi-centrale}

\Proved{} \textbf{Tesi.} Il destino del sistema globale nel 2031 sarà
deciso dalla distribuzione della \textbf{coerência de fase \(\Phi\)} tra
quattro attori strutturalmente diversi:

\begin{enumerate}
\def\labelenumi{\arabic{enumi}.}
\tightlist
\item
  \textbf{Soliton (Musk)} --- attore singolo con \(\Phi\) degenere ma
  sorgente esogena \(\zeta_{\text{Sol}}(t)\) di decoerência per i sistemi
  macro;
\item
  \textbf{Hegemon Decoerente (USA)} --- capacidade material residua ma
  \(\Phi \to 0\) per polarizzazione, hiperextensão, perdita di
  legitimidade esterna;
\item
  \textbf{Rígido Sincronizado (China)} --- \(\Phi^{\text{osservato}}\)
  massima per imposizione gerarchica, \(\Phi^{\text{spontaneo}}\) in
  colapso estrutural (Neijuan, demografia, decoupling);
\item
  \textbf{Balança Quântica (Brasil)} --- unico attore con
  condizioni strutturali per produrre \(\Phi\) esportabile,
  \emph{condizionalmente} all'esito elettorale del 4 ottobre 2026.
\end{enumerate}

\Hypoth{} \textbf{Risultato estrutural principale.} Solo il Bilanciere
produce \emph{\(\Phi\) esportabile vera}. Soliton, Egemone e Rigido
producono rispettivamente \(\zeta\) esogeno, decoerência distribuita, e
\(K\) imposto --- nessuno dei tre è \(\Phi\). La differenza è formale,
non retorica: emerge dalla struttura matematica di
\(T_s^{(2),\text{sys}}\) derivata nel Cap.~4 e consolidata nel Cap.~5.



\section{I Risultati Strutturali
Acquisiti}\label{i-risultati-strutturali-acquisiti}

\Proved{} \textbf{Risultati formali della Parte II (Cap.~3-5).}

{\def\LTcaptype{none} % do not increment counter
{\small\begin{longtable}[]{@{}
  >{\raggedright\arraybackslash}p{(\linewidth - 4\tabcolsep) * \real{0.3333}}
  >{\raggedright\arraybackslash}p{(\linewidth - 4\tabcolsep) * \real{0.3333}}
  >{\raggedright\arraybackslash}p{(\linewidth - 4\tabcolsep) * \real{0.3333}}@{}}
\toprule\noalign{}
\begin{minipage}[b]{\linewidth}\raggedright
\#
\end{minipage} & \begin{minipage}[b]{\linewidth}\raggedright
Risultato
\end{minipage} & \begin{minipage}[b]{\linewidth}\raggedright
Capítulo
\end{minipage} \\
\midrule\noalign{}
\endhead
\bottomrule\noalign{}
\endlastfoot
R-I & Identità \(\delta^* = 0{,}50\) tra Teorema Folk Discursivo
(Topografia) e Teorema Folk Geometrico (Geometria) --- invarianza di
scala dimostrata & Cap.~3.1.3 \\
R-II & Operatore di trasposizione
\(\langle X \rangle_\Phi := \Phi \cdot N^{-1}\sum_j X_j\) via Kuramoto &
Cap.~3.1bis \\
R-III & Forma chiusa
\(T_s^{(2),\text{sys}} = s_{iii} z \delta (1+\Phi\langle A\rangle) / [\sigma \varepsilon_{\text{dis}} (1 + \langle\Omega\rangle/\Phi)]\)
& Cap.~4.4 \\
R-IV & Asimmetria \(\Phi\)-moltiplica-\(A\) vs
\(\Phi\)-divide-\(\Omega\) --- firma della seconda legge termodinamica
nei sistemi sociali & Cap.~4.3.2 \\
R-V & Spazio estrutural \(\mathbb{V}^{(4)}\) con sei relazioni di
acoplamento bivariate R1-R6 & Cap.~5.1 \\
R-VI & Forma esplicita \(\Phi^*(A, \Omega, \varepsilon_{\text{dis}})\)
della superficie critica \(\mathcal{S}\) & Cap.~5.3 \\
R-VII & Sistema dinamico postulato
\(\dot\Phi, \dot A, \dot\Omega, \dot\varepsilon\) con coefficienti
realizzanti R1-R6 & Cap.~5.4 \\
\end{longtable}}
}

\Proved{} \textbf{Risultati applicativi della Parte III (Cap.~6-9).}

{\def\LTcaptype{none} % do not increment counter
{\small\begin{longtable}[]{@{}
  >{\raggedright\arraybackslash}p{(\linewidth - 4\tabcolsep) * \real{0.3333}}
  >{\raggedright\arraybackslash}p{(\linewidth - 4\tabcolsep) * \real{0.3333}}
  >{\raggedright\arraybackslash}p{(\linewidth - 4\tabcolsep) * \real{0.3333}}@{}}
\toprule\noalign{}
\begin{minipage}[b]{\linewidth}\raggedright
\#
\end{minipage} & \begin{minipage}[b]{\linewidth}\raggedright
Risultato
\end{minipage} & \begin{minipage}[b]{\linewidth}\raggedright
Capítulo
\end{minipage} \\
\midrule\noalign{}
\endhead
\bottomrule\noalign{}
\endlastfoot
R-VIII & Profilo discorsivo Musk \((8, 9, 2)\) = equilibrio separante
bayesiano della Topografia Ap.~A & Cap.~6.2 \\
R-IX & Soliton come sorgente \(\zeta(t)\) esogena nel sistema USA,
verificaçãota empiricamente 2025-2026 & Cap.~6.6 + Cap.~7.4 \\
R-X & Solitonizzazione discorsiva degli USA: \((7,6,8) \to (8,9,2)\),
firma su Brand Finance Soft Power 2026 & Cap.~7.3 \\
R-XI & Mappatura Neijuan (Xiang Biao) \(\to \Omega\) DEQ$^2$;
anti-involution campaign chinês come ammissione istituzionale & Cap.
8.2 \\
R-XII & Estado Logístico (Cervo) come \(\Phi^{\text{esportabile}}\)
operativa: sovrapposizione coerente del campo \(\Pi\) & Cap.~9.1 \\
R-XIII & Asimmetria triplice USA/China/Brasil: \(\Phi \to 0\) per
percorsi opposti, solo Brasil produce \(\Phi\) vera & Cap.~9.7 \\
\end{longtable}}
}

\Proved{} \textbf{Risultati operativi della Parte IV (Cap.~10-11).}

{\def\LTcaptype{none} % do not increment counter
{\small\begin{longtable}[]{@{}
  >{\raggedright\arraybackslash}p{(\linewidth - 4\tabcolsep) * \real{0.3333}}
  >{\raggedright\arraybackslash}p{(\linewidth - 4\tabcolsep) * \real{0.3333}}
  >{\raggedright\arraybackslash}p{(\linewidth - 4\tabcolsep) * \real{0.3333}}@{}}
\toprule\noalign{}
\begin{minipage}[b]{\linewidth}\raggedright
\#
\end{minipage} & \begin{minipage}[b]{\linewidth}\raggedright
Risultato
\end{minipage} & \begin{minipage}[b]{\linewidth}\raggedright
Capítulo
\end{minipage} \\
\midrule\noalign{}
\endhead
\bottomrule\noalign{}
\endlastfoot
R-XIV & Dashboard SPC-DEQ con sei pannelli + aggregato + cinque regole
WE-DEQ$^2$ + indice \(C_{\text{DEQ}}^2\) & Cap.~10 \\
R-XV & Quattro mondi possibili al 2031 (M1-M4) con probabilità
soggettive esplicite & Cap.~10.4 \\
R-XVI &
\(\Pr(\text{decoerenza globale 2031}) \approx 0{,}80\)-\(0{,}85\),
\(\Pr(\Phi\text{-engine}) \approx 0{,}55\)-\(0{,}70\) & Cap.~10.4 \\
R-XVII & 16 predições arriscadas datate con campi popperiani (1)-(5) &
Cap.~11 \\
R-XVIII & Classificazione: 5 strutturali, 8 locali, 2 meta.
\(\Pr(\text{fall. simult. strutt.}) \approx 0{,}018\) & Cap.~11.5 \\
\end{longtable}}
}

\Proved{} \textbf{Risultati normativi della Parte V (Cap.~12).}

{\def\LTcaptype{none} % do not increment counter
{\small\begin{longtable}[]{@{}
  >{\raggedright\arraybackslash}p{(\linewidth - 4\tabcolsep) * \real{0.3333}}
  >{\raggedright\arraybackslash}p{(\linewidth - 4\tabcolsep) * \real{0.3333}}
  >{\raggedright\arraybackslash}p{(\linewidth - 4\tabcolsep) * \real{0.3333}}@{}}
\toprule\noalign{}
\begin{minipage}[b]{\linewidth}\raggedright
\#
\end{minipage} & \begin{minipage}[b]{\linewidth}\raggedright
Risultato
\end{minipage} & \begin{minipage}[b]{\linewidth}\raggedright
Capítulo
\end{minipage} \\
\midrule\noalign{}
\endhead
\bottomrule\noalign{}
\endlastfoot
R-XIX & Matrice 3$\times$3 di protocolli operativi (livelli $\times$ orizzonti) con
condizione antifrágil esplicita & Cap.~12.1-12.4 \\
R-XX & Strategia \(S_{\mathcal{R}_-}\) per cittadini in sistemi sotto
soglia & Cap.~12.6 \\
R-XXI & Distinzione politica antifrágil vs realista/utopica come
superamento concettuale & Cap.~12.7 \\
\end{longtable}}
}

\Hypoth{} \textbf{Sintesi.} Ventuno risultati strutturali acquisiti,
organizzati in quattro parti, con tracciabilità formale a Topografia,
Geometria, MQNU. La DEQ$^2$ è ora un \emph{corpus teorico chiuso} ---
pronto per estensione (paper EN, articolo PT) e per falsificação (Cap.
11 protocollo).



\section{Posizione del Libro nella
Tradizione}\label{posizione-del-libro-nella-tradizione}

\Hypoth{} \textbf{Mappa esplicita.}

{\def\LTcaptype{none} % do not increment counter
{\small\begin{longtable}[]{@{}
  >{\raggedright\arraybackslash}p{(\linewidth - 6\tabcolsep) * \real{0.2500}}
  >{\raggedright\arraybackslash}p{(\linewidth - 6\tabcolsep) * \real{0.2500}}
  >{\raggedright\arraybackslash}p{(\linewidth - 6\tabcolsep) * \real{0.2500}}
  >{\raggedright\arraybackslash}p{(\linewidth - 6\tabcolsep) * \real{0.2500}}@{}}
\toprule\noalign{}
\begin{minipage}[b]{\linewidth}\raggedright
Tradizione
\end{minipage} & \begin{minipage}[b]{\linewidth}\raggedright
Eredità accolta
\end{minipage} & \begin{minipage}[b]{\linewidth}\raggedright
Mossa DEQ$^2$
\end{minipage} & \begin{minipage}[b]{\linewidth}\raggedright
Distanza dichiarata
\end{minipage} \\
\midrule\noalign{}
\endhead
\bottomrule\noalign{}
\endlastfoot
\textbf{Arrighi} --- cicli egemonici & finanziarizzazione terminale =
segnale di transizione & \(T_s^{(2),\text{sys}}\) formalizza la
transizione & la transizione 2028-2031 è quella di Arrighi
\emph{misurata} \\
\textbf{Turchin} --- structural-demographic & elite overproduction
\(\subset \Omega\); PSI \(\subset \varepsilon_{\text{dis}}\) & aggiunta
di \(\Phi\) come parâmetro mancante & Turchin cattura il denominatore,
manca il numeratore \\
\textbf{Wallerstein} --- world-system & Brasil non
\emph{semi-periferia} ma \(\Phi\)-engine & riposizionamento brasileiro
nello schema centro-periferia & rifiuto della linearità
centro-periferia \\
\textbf{Cervo} --- Estado Logístico & base operativa del Cap.~9 &
formalizzazione matematica via \(\Pi^{\text{Logístico}}\)
multi-vettoriale & radicamento esplicito \\
\textbf{Taleb} --- antifragilidade & \(A\) come
\(\partial^2\Pi/\partial\sigma^2\) & inserimento in apparato sistemico,
asimmetria \(\Phi\)-\(A\)/\(\Omega\) & da euristica a struttura \\
\textbf{Gramsci/Cox} --- egemonia critica & egemonia come operatore di
proiezione (eredità MQNU) & inserimento in DEQ$^2$ come \(s_{iii}\)
proiettata sull'asse \(z\) & continuità filosofica \\
\textbf{Luttwak} --- grande strategia & modello di registro: formalismo
+ compattezza & integrazione: registro tecnico + apparato matematico &
continuità di stile \\
\textbf{Escola de Frankfurt} (Marcuse, Adorno, Habermas) --- via
Topografia & microfundação comunicativa di \(\Pi\) & \(T_s^{(2)}\)
emerge dal campo \(\Pi\) topografico & la critica francofortese è ora
\emph{mensurável} \\
\textbf{Xiang Biao} --- Neijuan etnografato & \(\Omega\) DEQ$^2$ come
formalizzazione del Neijuan & da observação antropologica a variável
sistemica & formalizzazione \\
\textbf{Kuramoto} --- sincronização fisica & \(\Phi\) parâmetro
d'ordine & transposição de escala micro$\to$macro per accoppiamenti
emittente-ricevente & applicazione metodologica \\
\end{longtable}}
}

\Proved{} \textbf{Conseguenza posizionale.} La DEQ$^2$ \emph{non
sostituisce} nessuna delle tradizioni elencate --- le \emph{integra} in
un apparato in cui ciascuna riceve una collocazione formale. Il rifiuto
del manicheismo (DEQ$^2$ \emph{contro} X) è esplicito: ogni tradizione
cattura una parte vera e parziale dello stesso fenomeno.



\section{Agenda di Estensione}\label{agenda-di-estensione}

\Hypoth{} \textbf{Tre vettori già aperti.}

\subsection{\texorpdfstring{Paper EN --- \emph{From Discursive
Topography to Hegemonic Phase Coherence} (target:
\emph{Complexity})}{Paper EN --- From Discursive Topography to Hegemonic Phase Coherence (target: Complexity)}}\label{paper-en-from-discursive-topography-to-hegemonic-phase-coherence-target-complexity}

Stato attuale: scaffold completo (PAPER-EN-Topography-to-Hegemony.md,
2026-04-25). Lavoro residuo:

\begin{itemize}
\tightlist
\item
  Dataset 50+ paesi $\times$ 20+ anni (1995-2025) per calibrazione
  \(\Sigma^{\text{geo}} \leftarrow \Sigma^{\text{comm}}\) (Apêndice B);
\item
  Bootstrap \(n \geq 5000\) per intervalli di confidenza dei
  coefficienti del sistema dinamico Cap.~5.4;
\item
  Repository GitHub \texttt{bilanciere-quantistico/predizioni-deq2}
  operativo (Cap.~11.4);
\item
  Sottomissione \emph{Complexity} o, in alternativa, \emph{Physica A},
  \emph{EPJ B}, \emph{Journal of Conflict Resolution}.
\end{itemize}

\subsection{\texorpdfstring{Articolo PT --- versione brasileira (target:
\emph{Revista Brasilira de Política Internacional} o \emph{Contexto
Internacional})}{Articolo PT --- versione brasileira (target: Revista Brasilira de Política Internacional o Contexto Internacional)}}\label{articolo-pt-versione-brasileira-target-revista-brasileira-de-poluxedtica-internacional-o-contexto-internacional}

Focus: il Brasil come \(\Phi\)-engine condicional, l'\emph{Estado
Logístico} come politica formalizzabile, le predizioni B1-B4 con
dettaglio sui contesti regionali (Mercosul, BRICS, COP30 follow-up). Da
scrivere in portoghese, registro accademico brasileiro. Lunghezza target
\(20\)-\(25\) pagine. Tempistiche: post-elettorali (novembre 2026) per
integrare l'esito B1.

\subsection{Saggio IT/EN/ES/PT --- distribuzione
KDP}\label{saggio-itenespt-distribuzione-kdp}

Versione attuale: bozza completa di tutti i 13 capitoli + Apêndice B.
Lavoro residuo:

\begin{itemize}
\tightlist
\item
  Appendici A (derivazione formale completa), C (protocolli di
  medição --- preparata da \S{}5.6 e Cap.~10), D (dialogo testuale con
  tradizioni);
\item
  Revisione editoriale, illustrazioni concettuali per i sei pannelli
  SPC-DEQ;
\item
  Traduzioni EN/ES/PT;
\item
  Upload KDP via pipeline esistente.
\end{itemize}



\section{Tre Vettori Ancora da
Aprire}\label{tre-vettori-ancora-da-aprire}

\Hypoth{} \textbf{R\&D futuro non ancora avviato.}

\begin{enumerate}
\def\labelenumi{\arabic{enumi}.}
\tightlist
\item
  \textbf{Calibrazione real-time della Dashboard SPC-DEQ} come \emph{web
  service pubblico}: aggregare V-Dem, IMF, NBS, Pew, Brand Finance in
  cruscotto live aggiornato trimestralmente. Costo stimato:
  \(40\)-\(80\)k\$ + 1-2 anni di sviluppo. Possibile partnership con
  think tank (Brookings, Chatham House, Stimson, Itamaraty).
\item
  \textbf{Estensione del modello a sotto-attori non-statali}: applicare
  \(\boldsymbol{p}(t) \in \mathbb{V}^{(4)}\) a aziende, città, ONG
  transnazionali, comunità religiose. Necessita ricalibrazione delle sei
  relazioni R1-R6 sulle scale appropriate (la R4 \(A\)-\(\Omega\) è
  probabilmente meno forte a scala aziendale, onde la concorrenza spinge
  \(A\) e \(\Omega\) verso correlazione meno negativa).
\item
  \textbf{Integrazione formale con MQNU}: \(\Phi\) DEQ$^2$ come
  \emph{modulo} di un'ampiezza di probabilità nello spazio degli stati
  di MQNU? La connessione è sospettata ma non formalizzata. Potrebbe
  richiedere un terzo libro.
\end{enumerate}



\section{Inviti Espliciti alla
Falsificazione}\label{inviti-espliciti-alla-falsificação}

\Hypoth{} \textbf{Sette modi diretti per attaccare la DEQ$^2$.}

{\def\LTcaptype{none} % do not increment counter
{\small\begin{longtable}[]{@{}
  >{\raggedright\arraybackslash}p{(\linewidth - 4\tabcolsep) * \real{0.3333}}
  >{\raggedright\arraybackslash}p{(\linewidth - 4\tabcolsep) * \real{0.3333}}
  >{\raggedright\arraybackslash}p{(\linewidth - 4\tabcolsep) * \real{0.3333}}@{}}
\toprule\noalign{}
\begin{minipage}[b]{\linewidth}\raggedright
\#
\end{minipage} & \begin{minipage}[b]{\linewidth}\raggedright
Critica
\end{minipage} & \begin{minipage}[b]{\linewidth}\raggedright
Risposta DEQ$^2$ (e onde cercare)
\end{minipage} \\
\midrule\noalign{}
\endhead
\bottomrule\noalign{}
\endlastfoot
1 & L'asimmetria \(\Phi\)-moltiplica-\(A\) vs \(\Phi\)-divide-\(\Omega\)
è arbitraria & Cap.~4.3.2: argomento termodinamico; alternative non
lineari da testare nel paper EN \\
2 & \(\Phi\) di Kuramoto non è applicabile a sistemi sociali & Cap.~4.1:
hipótese di campo médio è esplicita; alternative (correlazione di
Pearson, mutual information, Granger causality) producono risultati
equivalenti per sistemi \(N\) grande \\
3 & Le 16 predizioni Cap.~11 sono troppo numerose per essere
genuinamente rischiose & Cap.~11.5: 5 sono \emph{strutturali}, ciascuna
individualmente capace di invalidare l'apparato \\
4 & La centralità di B1 (elezione brasileira) è interessata & Cap.~9
\S{}9.0 + Profilo Utente: l'autore ha dichiarato esplicitamente la propria
origine brasileira e ha mantenuto onestà analitica con tre fragilidade
esplicitate \\
5 & Il modello sottostima eventi estremi (cigni neri) & Cap.~12.6:
\(S_{\mathcal{R}_-}\) è progettato esplicitamente per scenari peggiori
del \emph{worst} mediano \\
6 & L'apparato Estado Logístico è specificamente brasileiro e non
generalizzabile & Cap.~9.1: la struttura formale di
\(\Pi^{\text{Logístico}}\) multi-vettoriale è generalizzabile a
\emph{qualunque} attore con \(s_{iii} > s_{ii}\) relativo --- Brasil è
il caso paradigmatico, non l'unico possibile \\
7 & La DEQ$^2$ ignora variáveis climatiche/ecologiche & Cap.~4.1 + Cap.~5:
\(\sigma\) esogena include shock climatici; COP30 è incorporata nel Cap.
9.2 come test di \(\Phi^{\text{esportabile}}\). Estensione formale
dell'asse climatico è agenda futura \\
\end{longtable}}
}

\Proved{} \textbf{Tutte le critiche sopra producono \emph{aggiornamento}
della teoria, non sua sospensione}. La DEQ$^2$ è progettata per crescere
via critica.



\section{Chiusura}\label{chiusura}

\Hypoth{} \textbf{Ultimo paragrafo del libro.}

Le predizioni di questo libro saranno verificaçãote o falsificate fra
ottobre 2026 (B1 --- elezione brasileira) e dicembre 2031 (G1, G2 ---
sistema globale). Nel frattempo il sistema globale opera, decide,
sceglie. La DEQ$^2$ non pretende di prescrivere quelle scelte --- pretende
di \emph{renderle leggibili}, di mostrare il loro spazio degli effetti,
di rendere trasparente il costo dell'inazione e l'asimmetria del
rischio. Se il suo apparato si dimostrerà robusto, sarà porque è stato
attaccato; se verrà superato, sarà da una teoria che fa meglio le stesse
cose. In entrambi i casi avrà fatto il suo lavoro.

Il sistema globale entra nel triennio 2028-2031 con coerência de fase in
declino quasi ovunque, capacidade material concentrata in attori
strutturalmente fragili, e una sola candidatura non degenere a
\(\Phi\)-engine globale. Quella candidatura sarà confermata o smentita
dal Brasil il 4 ottobre 2026. Da quel giorno in depois, ogni cittadino,
ogni organizzazione, ogni Stato avrà un dato in più per scegliere fra
strategie passive e strategie antifragili. Questo libro è scritto porque
quella scelta sia \emph{informata}, non porque sia
\emph{predeterminata}.



\section{Notas Epistêmicas de Encerramento del
Libro}\label{note-epistemiche-di-chiusura-del-libro}

{\def\LTcaptype{none} % do not increment counter
{\small\begin{longtable}[]{@{}
  >{\raggedright\arraybackslash}p{(\linewidth - 4\tabcolsep) * \real{0.3333}}
  >{\raggedright\arraybackslash}p{(\linewidth - 4\tabcolsep) * \real{0.3333}}
  >{\raggedright\arraybackslash}p{(\linewidth - 4\tabcolsep) * \real{0.3333}}@{}}
\toprule\noalign{}
\begin{minipage}[b]{\linewidth}\raggedright
\#
\end{minipage} & \begin{minipage}[b]{\linewidth}\raggedright
Afirmação
\end{minipage} & \begin{minipage}[b]{\linewidth}\raggedright
Status
\end{minipage} \\
\midrule\noalign{}
\endhead
\bottomrule\noalign{}
\endlastfoot
1 & Tesi centrale ricapitolata fedelmente & \Proved{} \\
2 & 21 risultati strutturali tracciati ai capitoli di origine &
\Proved{} aggregazione completa \\
3 & Posizione esplicita verso 10 tradizioni & \Hypoth{} mappa
concettuale \\
4 & Tre vettori di estensione già aperti (paper EN, articolo PT, saggio
multi-lingua) & \Proved{} pianificati \\
5 & Tre vettori R\&D futuri (web service, sub-attori non-statali,
integrazione MQNU) & \Conj{} prospettici \\
6 & Sette inviti espliciti alla falsificação con risposte DEQ$^2$ &
\Hypoth{} esposizione metodologica \\
7 & Chiusura politica: la DEQ$^2$ come rendering visibile, non come
prescrizione & \Hypoth{} posizione esplicita \\
\end{longtable}}
}



\section{Ringraziamenti}\label{ringraziamenti}

A \textbf{Clayton Teixeira}, collaboratore principale del corpus
matematico (em particular \emph{La Costante di Marcelino}, base
epistemica dell'apparato formale). Alla \textbf{tradizione critica
francofortese} che ha reso pensabile la formalizzazione della
comunicazione come oggetto matematico. Ai \textbf{lettori che vorranno
falsificare} anziché accogliere --- il loro lavoro è la condizione di
vita della teoria.


